\section{Discussion}
\label{sec:discussion}

\subsection{What the measurements show}
With patient- and duplicate-grouped splits and no proxy labels, a standard DenseNet-121 reached an internal test accuracy of \RTestAcc\% and a macro ROC-AUC of \RTestMacroAUC. On radiographs from a hospital it had never seen, it flagged \RExtSensK/\RExtSensN{} TB scans as abnormal but called only \RExtSpecK/\RExtSpecN{} normal scans Normal, and named \RExtTBK/\RExtTBN{} TB scans Tuberculosis. The internal results therefore overstate behaviour at an unseen hospital by a wide margin, although the usual leakage paths were closed for sources with patient identifiers. Closing the leakage paths made the internal estimate credible but left it a poor guide to generalisation.

The new analyses address one possible reason, source confounding. Source information is available throughout the spectrum: the source of an image can be recovered from every one of six octave bands, including within Normal images, and removing a single band changes little (Section~\ref{sec:res-c1}). The information is also entangled with disease information: twenty steps of linear erasure lowered the source accuracy a fresh probe could reach from \RInlpImgTestSrcBAKZero\% to \RInlpImgTestSrcBAKStar\%, not to chance, and cost disease performance, and erasing the linearly decodable information in one step cost more (Section~\ref{sec:res-c3}). Finally, a linear head trained without a source scored far below its in-distribution reference on four of five held-out sources (Section~\ref{sec:res-c4}). These are measurements on frozen features. Together with the per-source and external results they are consistent with a model that depends on source-specific cues, but each is indirect: a probe shows that information is available, not that the trained network uses it, and we make no causal claim about the shortcut or about its location. For the same reason we do not write that source information ``cannot be removed'': only linear erasure, a 20-iteration cap, one cache, one benchmark and one run were examined, and a probe near chance would not prove absence either.

\paragraph{Novelty and its limits} To our knowledge, within the searches described, no prior work quantifies, band by band, source versus disease information in CNN features of chest radiographs. That statement concerns only the localisation in C1, and its main outcome is that the signal is not localised. Per-band harmonisation~\cite{kim2023homogenization} and spectral failure analysis~\cite{halimi2026spectral} exist for chest radiographs, so we do not claim frequency-filtered source classification as new in general. Linear erasure~\cite{ravfogel2020null,belrose2023leace} and leave-one-source-out evaluation~\cite{yao2019strong} are published; here they are measurements, and neither is claimed as a method or a protocol contribution. The negative C2 decision is part of the contribution: the rule, fixed in advance, prevented a retraining that the frozen-feature evidence did not support, and we make no claim about band suppression in a trained network.

\paragraph{Caveats that bound these readings}
(i)~The cache is resampled to $224\times224$ after filtering, and sources differ in native resolution and field of view; the high bands in particular may reflect resampling and the pre-filters, and the lower high-band source accuracy of the no-filter cache is consistent with that. No claim about scanner physics is made. (ii)~The secondary extractor was trained on all six sources, so trained-backbone numbers are not held-out. (iii)~Disease and source are confounded (Effusion and Nodule/Mass occur only in NIH), which is why the within-source disease score is primary. (iv)~Intervals cover test-sample variation only, with one probe fit per variant; samples are small (150 and 100 images per cell) and only five sources have two classes with at least 20 images. (v)~The NIH gap in C4 is a metric artefact (Section~\ref{sec:res-c4}). (vi)~Patient identity is unknown for TBX11K and COVID-19 Radiography rows and derived from file names for Kermany, so every held-out number carries that caveat. (vii)~Montgomery is not untouched: it was scored for runs 1--5 and in C3 and C4, and run~2 was started after run~1's external results were seen. (viii)~Everything is one run per analysis with no seeds.

\subsection{Ranking versus operating point}
The external ROC-AUC of the abnormality score (\RExtAUC) is well above chance while specificity is low, just as internal macro ROC-AUC is much higher than internal accuracy. A shift of scores toward ``abnormal'' at a new hospital can leave the ranking partly intact and still move most normal images across the arg-max boundary. A model deployed at a new site would need its operating point re-established there. We do not compare these numbers with those of other studies, which used different data and splits.

\subsection{What the audit of the earlier pipeline showed}
\label{sec:disc-audit}
The V5 pipeline pooled eight public sources, several from third-party redistributions, and recorded an internal test accuracy of \RVFiveAcc\% on $N=\RVFiveN$ images. A retrospective audit found that 154 of the 1,570 internal test images (9.8\%) belonged to patients who also had images in the training split, and that 8,478 further images had no real patient identity, so their separation could not be checked. About half of the Pneumonia class (1,395 labels) consisted of TBX11K images labelled ``sick but non-TB'', 124 Tuberculosis labels were latent TB, and 1,440 of 1,467 VinDr-CXR images used as single-label examples carried two or more findings. On Montgomery every recorded V5 variant flagged essentially every TB radiograph as abnormal but called none of the normal radiographs Normal (specificity \RVFiveExtSpecK/\RVFiveExtSpecN); a sensitivity quoted alone would have suggested a perfect screening tool. One V5 TB source is described as containing Montgomery images, and a content check could neither confirm nor exclude the overlap. Earlier drafts also contained reporting errors, and the canonical checkpoint was lost. V6 removes each defect or measures it explicitly. The V5 numbers are not comparable with V6.

\subsection{Implications for multi-source benchmark design}
The audit and the analyses suggest the following practices: acquire from official sources and log provenance per file; fix labelling rules before training and label only from the source's own labels; split by patient and by content, stating how many images have unknown identity; report results per source and flag classes supplied by one source; evaluate on an unseen hospital and always report specificity with sensitivity; and fix selection rules in advance (here recorded in the repository, not registered externally). Because source information is linearly available throughout the spectrum, these results give no reason to expect that band filtering or linear erasure of frozen features would make up for source confounding; source-balanced data were not tested.

\subsection{Input gate}
\label{sec:disc-gate}
The application refuses images judged not to be chest radiographs, using a heuristic score because the intended trained gate was unavailable. At the inherited threshold (0.83) it rejected \RCxrGateBOneOldTestRej{} of \RCxrGateBOneOldTestN{} genuine radiographs of the build-1 test split (\RCxrGateBOneOldTestRejPct\%). After recalibration on build-2 validation radiographs only (threshold \RCxrGateNewThr), it rejected \RCxrGateBTwoNewTestRej{} of \RCxrGateBTwoOldTestN{} build-2 test radiographs (\RCxrGateBTwoNewTestRejPct\%), against \RCxrGateBTwoOldTestRej{} (\RCxrGateBTwoOldTestRejPct\%) at 0.83 on the same split. The negatives used for calibration are synthetic and few, so the gate's ability to reject real non-radiographs is not established.
\subsection{Limitations}
\label{sec:limitations}
\begin{itemize}
\item \textbf{Frozen features, linear methods, one benchmark.} C1, C3 and C4 examine linear probes and heads on frozen DenseNet-121 features of one cache, with six sources and one run each. They say nothing about fine-tuned networks, nonlinear erasure or other data. No retraining per held-out source was done, so the served model is not shown to behave as the C4 heads do.
\item \textbf{Few training runs, CPU, time-boxed.} Run~1 is a single run, run~2 has two further seeds, and the ablation has one seed; the SD of test accuracy (\RMSTestAccSD{} points) is a rough estimate that does not extend to other splits, data or architectures.
\item \textbf{Limited ablation.} Only the full chain and no filtering were compared; single filters, the order relative to resizing and alternative backbones were not studied. Run~5 selected its last epoch.
\item \textbf{One small, binary external cohort}, not strictly untouched.
\item \textbf{Label quality and the single-label assumption.} NIH labels are mined from reports; images with several findings were excluded.
\item \textbf{Population and source confounding.} Kermany images are paediatric, all others adult; no demographic data were available, so subgroup performance~\cite{seyyedkalantari2021underdiagnosis} could not be assessed.
\item \textbf{Uncalibrated scores and unvalidated thresholds.} No calibration map was fitted, and the application's 0.5 ``inconclusive'' threshold is not validated.
\item \textbf{No clinical validation.} The system has not been evaluated prospectively or by radiologists and must not be used for diagnosis.
\end{itemize}

\subsection{Future work}
In order: a causal test of the source shortcut in a trained network (leave-one-source-out training and lung-masked inference); the C2 band-stop retraining only if new evidence justifies it, under a rule again fixed in advance; a phase-versus-amplitude split per band; nonlinear erasure; more seeds, so that runs can be compared as distributions; a second backbone; access to VinDr-CXR and JSRT so that Effusion and Nodule/Mass have more than one source; a second external cohort with more classes, ideally a local set collected with ethics approval; and evaluation of the input gate on real non-radiographs.
